\section{AHSME 1974}
        \begin{enumerate}
        	\item (\textit{AHSME 1974}) If $x\not=0$ or $4$ and $y\not=0$ or $6$, then $\frac{2}{x}+\frac{3}{y}=\frac{1}{2}$ is equivalent to


 $\mathrm{(A)\ } 4x+3y=xy \qquad \mathrm{(B) \ }y=\frac{4x}{6-y} \qquad \mathrm{(C) \  } \frac{x}{2}+\frac{y}{3}=2 \qquad$


 $\mathrm{(D) \  } \frac{4y}{y-6}=x \qquad \mathrm{(E) \  }\text{none of these}$


 
 \textbf{Answer: }D



\textbf{Solution}

To clear out the fractions, we multiply both sides of the equation by $xy$ to get $2y+3x=\frac{xy}{2}$, or $4y+6x=xy$. From this, we have $4y=xy-6x$. We can factor an $x$ out of the RHS, and so $4y=x(y-6)\implies \frac{4y}{y-6}=x, \boxed{\text{D}}$.

 


	\item (\textit{AHSME 1974}) Let $x_1$ and $x_2$ be such that $x_1\not=x_2$ and $3x_i^2-hx_i=b$, $i=1, 2$. Then $x_1+x_2$ equals


 $\mathrm{(A)\ } -\frac{h}{3} \qquad \mathrm{(B) \ }\frac{h}{3} \qquad \mathrm{(C) \  } \frac{b}{3} \qquad \mathrm{(D) \  } 2b \qquad \mathrm{(E) \  }-\frac{b}{3}$


 
 \textbf{Answer: }B



\textbf{Solution}

Notice that $x_1$ and $x_2$ are the distinct solutions to the quadratic $3x^2-hx-b=0$. By Vieta, the sum of the roots of this quadratic is the negation of the coefficient of the linear term divided by the coefficient of the quadratic term, so in this case $-\frac{-h}{3}=\frac{h}{3}, \boxed{\text{B}}$.

 


	\item (\textit{AHSME 1974}) The coefficient of $x^7$ in the polynomial expansion of 


 
 \[(1+2x-x^2)^4\]


 is


 $\mathrm{(A)\ } -8 \qquad \mathrm{(B) \ }12 \qquad \mathrm{(C) \  } 6 \qquad \mathrm{(D) \  } -12 \qquad \mathrm{(E) \  }\text{none of these}$


 
 \textbf{Answer: }A



\textbf{Solution 1}

Let's write out the multiplication, so that it becomes easier to see.

 
 \[(1+2x-x^2)(1+2x-x^2)(1+2x-x^2)(1+2x-x^2)\]

 We can now see that the only way to get an $x^7$ is by taking three $-x^2$ and one $2x$. There are $\binom{4}{1}=4$ ways to pick which term the $2x$ comes from, and the coefficient of each one is $(-1)^3(2)=-2$. Therefore, the coefficient of $x^7$ is $(4)(-2)=-8, \boxed{\text{A}}$.

 



\textbf{Solution 2}


 \[(x^2-2x-1)^4=(x^2-2x+1-2)^4=((x-1)^2-(\sqrt{2})^2)^4=(x-1-\sqrt{2})^4(x-1+\sqrt{2})^4\]Each of the two polynomials can have terms of order 4,3,2,1,0. The only way to obtain a term of order 7 is either a term of order 4 from the first and 3 from the second, or the other way around. Notice that the terms of order 4 must necessarily have coefficient 1. Also keeping in mind that the coefficient of $x^3$ in $(x+a)^4$ is 4a, we have:

 - Coefficient of $x^3$ (by extension $x^7$) in the first polynomial $=4\times (-1-\sqrt{2})$;

 - Coefficient of $x^3$ (by extension $x^7$) in the second polynomial is $4\times (-1+\sqrt{2})$.

 So the coefficient of $x^7$ in the expansion is $(-4-4\sqrt{2})+(-4+4\sqrt{2})=-4-4=-8, \boxed{\text{A}}$.

 


	\item (\textit{AHSME 1974}) What is the remainder when $x^{51}+51$ is divided by $x+1$?


 $\mathrm{(A)\ } 0 \qquad \mathrm{(B) \ }1 \qquad \mathrm{(C) \  } 49 \qquad \mathrm{(D) \  } 50 \qquad \mathrm{(E) \  }51$


 
 \textbf{Answer: }D



\textbf{Solution 1}

From the \href{https://artofproblemsolving.com/wiki/index.php?title=Remainder\_Theorem}{Remainder Theorem}, the remainder when $x^{51}+51$ is divided by $x+1$ is $(-1)^{51}+51=-1+51=50, \boxed{\text{D}}$.

 



\textbf{Solution 2}

Notice that $x^{51}+51$ is equal to $(x^{51}+1)+50$. Since $x^51+1=(x+1)(x^{50}-x^{49}+ \ldots +x^2 - x + 1)$, this part of the polynomial is divisible by $x+1$. Thus, the remainder is $\boxed{\textbf{(D)}~50}$.

 


	\item (\textit{AHSME 1974}) Given a quadrilateral $ABCD$ inscribed in a circle with side $AB$ extended beyond $B$ to point $E$, if $\measuredangle BAD=92^\circ$ and $\measuredangle ADC=68^\circ$, find $\measuredangle EBC$.


 $\mathrm{(A)\ } 66^\circ \qquad \mathrm{(B) \ }68^\circ \qquad \mathrm{(C) \  } 70^\circ \qquad \mathrm{(D) \  } 88^\circ \qquad \mathrm{(E) \  }92^\circ$


 
 \textbf{Answer: }B



\textbf{Solution}

Since $ABCD$ is cyclic, opposite angles must sum to $180^\circ$. Therefore, $\angle ADC+\angle ABC=180^\circ$, and $\angle ABC=180^\circ-\angle ADC=180^\circ-68^\circ=112^\circ$. Notice also that $\angle ABC$ and $\angle CBE$ form a linear pair, and so they sum to $180^\circ$. Therefore, $\angle EBC=180^\circ-\angle ABC=180^\circ-112^\circ=68^\circ, \boxed{\text{B}}$. Notice that the answer didn't even depend on $\angle BAD$.

 


	\item (\textit{AHSME 1974}) For positive real numbers $x$ and $y$ define $x*y=\frac{x\cdot y}{x+y}$' then


 $\mathrm{(A)\ } \text{``*" is commutative but not associative} \qquad$


 $\mathrm{(B) \ }\text{``*" is associative but not commutative} \qquad$


 $\mathrm{(C) \  } \text{``*" is neither commutative nor associative} \qquad$


 $\mathrm{(D) \  } \text{``*" is commutative and associative} \qquad$


 $\mathrm{(E) \  }\text{none of these} \qquad$


 
 \textbf{Answer: }D



\textbf{Solution}

First, let's check for commutivity. We have 
 \[y*x=\frac{y\cdot x}{y+x}=\frac{x\cdot y}{x+y}=x*y\], so $*$ is commutative.

 Now we check for associativity. We have 
 \[(x*y)*z=\left(\frac{x\cdot y}{x+y}\right)*z=\frac{\frac{x\cdot y}{x+y}\cdot z}{\frac{x\cdot y}{x+y}+z}=\frac{x\cdot y\cdot z}{x\cdot y+z(x+y)}=\frac{x\cdot y\cdot z}{x\cdot y+y\cdot z+x\cdot z}.\] Also, 
 \[x*(y*z)=x*\left(\frac{y\cdot z}{y+z}\right)=\frac{x\cdot\frac{y\cdot z}{y+z}}{x+\frac{y\cdot z}{y+z}}=\frac{x\cdot y\cdot z}{x(y+z)+y\cdot z}=\frac{x\cdot y\cdot z}{x\cdot y+y\cdot z+x\cdot z}=(x*y)*z ,\] and so $*$ is also associative. $\boxed{\text{D}}$

 


	\item (\textit{AHSME 1974}) A town's population increased by $1,200$ people, and then this new population decreased by $11\%$. The town now had $32$ less people than it did before the $1,200$ increase. What is the original population?


 $\mathrm{(A)\ } 1,200 \qquad \mathrm{(B) \ }11,200 \qquad \mathrm{(C) \  } 9,968 \qquad \mathrm{(D) \  } 10,000 \qquad \mathrm{(E) \  }\text{none of these}$


 
 \textbf{Answer: }D



\textbf{Solution}

Let $n$ be the original population. When $1,200$ people came, the new population was $n+1200$. Then, since the new population decreased by $11\%$, we must multiply this new population by $.89$ to get the final population. This is equal to $(0.89)(n+1200)=0.89n+1068$. We're given that this is $32$ less than the original popluation, $n$. Therefore, $0.89n+1068=n-32\implies 0.11n=1100\implies n=\frac{1100}{0.11}=10000, \boxed{\text{D}}$.

 


	\item (\textit{AHSME 1974}) What is the smallest prime number dividing the sum $3^{11}+5^{13}$?


 $\mathrm{(A)\ } 2 \qquad \mathrm{(B) \ }3 \qquad \mathrm{(C) \  } 5 \qquad \mathrm{(D) \  } 3^{11}+5^{13} \qquad \mathrm{(E) \  }\text{none of these}$


 
 \textbf{Answer: }A



\textbf{Solution}

Since we want to find the smallest prime dividing the sum, we start with the smallest prime and move up, so first we try $2$. Notice that $3^{11}$ and $5^{13}$ are both odd, so their sum must be even. This means that $2$ must divide $3^{11}+5^{13}$, and so since $2$ is the smallest prime, our answer must be $2, \boxed{\text{A}}$.

 


	\item (\textit{AHSME 1974}) The integers greater than one are arranged in five columns as follows:


 
 \[\begin{tabular}{c c c c c}\ & 2 & 3 & 4 & 5\\ 9 & 8 & 7 & 6 &\ \\ \ & 10 & 11 & 12 & 13\\ 17 & 16 & 15 & 14 &\ \\ \ & . & . & . & .\\ \end{tabular}\]


 (Four consecutive integers appear in each row; in the first, third and other odd numbered rows, the integers appear in the last four columns and increase from left to right; in the second, fourth and other even numbered rows, the integers appear in the first four columns and increase from right to left.)


 In which column will the number $1,000$ fall?


 $\mathrm{(A)\ } \text{first} \qquad \mathrm{(B) \ }\text{second} \qquad \mathrm{(C) \  } \text{third} \qquad \mathrm{(D) \  } \text{fourth} \qquad \mathrm{(E) \  }\text{fifth}$


 
 \textbf{Answer: }B



\textbf{Solution}

We try pairing numbers with the column number they're in.

 
 \[\begin{tabular}{c c}\ \text{Number} & \text{Column Number} \\ 2 & 2\\ 3 & 3\\ 4 & 4\\ 5 & 5\\ 6 & 4\\ 7 & 3\\ 8 & 2\\ 9 & 1\\ 10 & 2\\ 11 & 3\\ 12 & 4\\ \end{tabular}\]

 Now we can see a pattern. The column number starts at $2$, goes up to $5$, then goes back down to $1$, and repeats. Each of these blocks has a length of $7$, so the column number has a period of $8$. Since $1000\equiv8(\bmod\,8)$, $1000$ is in the same column as $8$, which is the second column. $\boxed{\text{B}}$

 


	\item (\textit{AHSME 1974}) What is the smallest integral value of $k$ such that 


 
 \[2x(kx-4)-x^2+6=0\]


 has no real roots?


 $\mathrm{(A)\ } -1 \qquad \mathrm{(B) \ }2 \qquad \mathrm{(C) \  } 3 \qquad \mathrm{(D) \  } 4 \qquad \mathrm{(E) \  }5$


 
 \textbf{Answer: }B



\textbf{Solution}

Expanding, we have $2kx^2-8x-x^2+6=0$, or $(2k-1)x^2-8x+6=0$. For this quadratic not to have real roots, it must have a negative discriminant. Therefore, $(-8)^2-4(2k-1)(6)<0\implies 64-48k+24<0\implies k>\frac{11}{6}$. From here, we can easily see that the smallest integral value of $k$ is $2, \boxed{\text{B}}$.

 


	\item (\textit{AHSME 1974}) If $(a, b)$ and $(c, d)$ are two points on the line whose equation is $y=mx+k$, then the distance between $(a, b)$ and $(c, d)$, in terms of $a, c,$ and $m$ is 


 $\mathrm{(A)\ } |a-c|\sqrt{1+m^2} \qquad \mathrm{(B) \ }|a+c|\sqrt{1+m^2} \qquad \mathrm{(C) \  } \frac{|a-c|}{\sqrt{1+m^2}} \qquad$


 $\mathrm{(D) \  } |a-c|(1+m^2) \qquad \mathrm{(E) \  }|a-c|\,|m|$


 
 \textbf{Answer: }A



\textbf{Solution}

Notice that since $(a, b)$ is on $y=mx+k$, we have $b=am+k$. Similarly, $d=cm+k$. Using the distance formula, the distance between the points $(a, b)$ and $(c, d)$ is

 $\sqrt{(a-c)^2+(b-d)^2}=\sqrt{(a-c)^2+(am+k-cm-k)^2}$ 

 $=\sqrt{(a-c)^2+m^2(a-c)^2}$ 

 $=|a-c|\sqrt{1+m^2}.$

 And so the answer is $\boxed{\text{A}}$.

 


	\item (\textit{AHSME 1974}) If $g(x)=1-x^2$ and $f(g(x))=\frac{1-x^2}{x^2}$ when $x\not=0$, then $f(1/2)$ equals 


 $\mathrm{(A)\ } 3/4 \qquad \mathrm{(B) \ }1 \qquad \mathrm{(C) \  } 3 \qquad \mathrm{(D) \  } \sqrt{2}/2 \qquad \mathrm{(E) \  }\sqrt{2}$


 
 \textbf{Answer: }B



\textbf{Solution}

Notice that if we can find a $y$ such that $g(y)=\frac{1}{2}$, then we can plug that into the second equation to get $f\left(\frac{1}{2}\right)=f(g(y))=\frac{1-y^2}{y^2}$. Thus, we let $1-y^2=\frac{1}{2}\implies y^2=\frac{1}{2}$. Notice that, since our final expression only involves $y^2$, we don't need to take the square root. Thus, $f\left(\frac{1}{2}\right)=\frac{1-y^2}{y^2}=\frac{1-\frac{1}{2}}{\frac{1}{2}}=1, \boxed{\text{B}}$.

 


	\item (\textit{AHSME 1974}) Which of the following is equivalent to "If P is true, then Q is false."?


 $\mathrm{(A)\ } \text{``P is true or Q is false."} \qquad$


 $\mathrm{(B) \ }\text{``If Q is false then P is true."} \qquad$


 $\mathrm{(C) \  } \text{``If P is false then Q is true."} \qquad$


 $\mathrm{(D) \  } \text{``If Q is true then P is false."} \qquad$


 $\mathrm{(E) \  }\text{``If Q is true then P is true."} \qquad$


 
 \textbf{Answer: }D



\textbf{Solution}

Remember that a statement is logically equivalent to its contrapositive, which is formed by first negating the hypothesis and conclusion and then switching them. In this case, the contrapositive of "If P is true, then Q is false." is "If Q is true, then P is false." $\boxed{\text{D}}$

 The fact that a statement's contrapositive is logically equivalent to it can easily be seen from a venn diagram arguement.

 
                    \begin{figure}[H]
                        centering
                        \includegraphics[width=0.5\textwidth]{Solution/AHSME_I_1974_13_Solution_Figure_1.png}
                    \end{figure}
                    

 From this venn diagram, clearly "If $p$, then $q$." is true. However, since $p$ is fully contained in $q$, the statement "If not $q$, then not $p$." is also true, and so a statement and its contrapositive are equivalent.

 



\textbf{Solution}

Let's consider the following case of dogs. If p = true then the dog is 100\% blue. If q = true then the dog is 100\% red. 

 This is a valid solution to if P is true, then Q is false, or if the dog is 100\% blue then the dog is not 100\% red.

 Now let's look at the options:$\mathrm{(A)\ } \text{``P is true or Q is false."} \qquad$ ---> If the dog is 100\% blue or the dog is not 100\% red. This statement makes no sense.

 $\mathrm{(B) \ }\text{``If Q is false then P is true."} \qquad$ ---> If the dog is not 100\% red, then the dog is 100\% blue. This does not make sense, as the dog could be other colors, like white, or black, or yellow. My dog is brown!

 $\mathrm{(C) \  } \text{``If P is false then Q is true."} \qquad$ ---> If the dog is not 100\% blue, then the dog is 100\% red. Again, this fails for other colors of dogs. This is dogscrimination!

 $\mathrm{(D) \  } \text{``If Q is true then P is false."} \qquad$ ---> If the dog is 100\% red, then it is not 100\% blue. This statement makes logical sense, and is correct.

 $\mathrm{(E) \  }\text{``If Q is true then P is true."} \qquad$ --->  If the dog is 100\% red, then the dog is 100\% blue. This does not make sense, because we know the color of the dog, a dog cannot be two colors at once. Actually in quantum mechanics maybe, but not for this AHSME problem.(This example is not the best way to test this answer choice if there were a different question so mind that by TFEA)

 The statement that makes the most sense is $\boxed{\text{D}}$

 


	\item (\textit{AHSME 1974}) Which statement is correct?


 $\mathrm{(A)\ } \text{If } x<0, \text{then } x^2>x. \qquad \mathrm{(B) \ } \text{If } x^2>0, \text{then } x>0.$


 $\qquad \mathrm{(C) \  } \text{If } x^2>x, \text{then } x>0. \qquad \mathrm{(D) \  } \text{If } x^2>x, \text{then } x<0.$


 $\qquad \mathrm{(E) \  }\text{If } x<1, \text{then } x^2<x.$


 
 \textbf{Answer: }A



\textbf{Solution}

We begin by going statement by statement.

 For (A), $x^2>0$ for all real $x$, and so $x^2>0>x$. Therefore, $\boxed{\text{A}}$ is true.

 We can give counterexamples for all other statements:

 For (B), $x=-1$

 For (C), $x=-1$

 For (D), $x=2$

 For (E), $x=-1$

 


	\item (\textit{AHSME 1974}) If $x<-2$, then $|1-|1+x||$ equals 


 $\mathrm{(A)\ } 2+x \qquad \mathrm{(B) \ }-2-x \qquad \mathrm{(C) \  } x \qquad \mathrm{(D) \  } -x \qquad \mathrm{(E) \  }-2$


 
 \textbf{Answer: }B



\textbf{Solution}

Notice that, for $a<0$, $|a|=-a$, and for $a>0$, $|a|=a$.

 Since $x<-2$, $1+x<1-2<0$, so $|1+x|=-x-1$. Therefore, $1-|1+x|=1+x+1=x+2<-2+2=0$, and so $|1-|1+x||=|x+2|=-2-x, \boxed{\text{B}}$.

 


	\item (\textit{AHSME 1974}) A circle of radius $r$ is inscribed in a right isosceles triangle, and a circle of radius $R$ is circumscribed about the triangle. Then $R/r$ equals


 $\mathrm{(A)\ } 1+\sqrt{2} \qquad \mathrm{(B) \ }\frac{2+\sqrt{2}}{2} \qquad \mathrm{(C) \  } \frac{\sqrt{2}-1}{2} \qquad \mathrm{(D) \  } \frac{1+\sqrt{2}}{2} \qquad \mathrm{(E) \  }2(2-\sqrt{2})$


 
 \textbf{Answer: }A



\textbf{Solution}


                    \begin{figure}[H]
                        centering
                        \includegraphics[width=0.5\textwidth]{Solution/AHSME_I_1974_16_Solution_Figure_1.png}
                    \end{figure}
                    

 Label the points as in the figure above. Let the side length $AB=AC=s$. Therefore, $BC=s\sqrt{2}$. Since the circumradius of a right triangle is equal to half of the length of the hypotenuse, we have $R=\frac{s\sqrt{2}}{2}$.

 Now to find the inradius. Notice that $IFAE$ is a square with side length $r$, and also $AD=R$. Therefore, $s=AD=AI+ID=r\sqrt{2}+r$, and so $r=\frac{R}{\sqrt{2}+1}$.

 Finally, $\frac{R}{r}=1+\sqrt{2}, \boxed{\text{A}}$.

 


	\item (\textit{AHSME 1974}) If $i^2=-1$, then $(1+i)^{20}-(1-i)^{20}$ equals


 $\mathrm{(A)\ } -1024 \qquad \mathrm{(B) \ }-1024i \qquad \mathrm{(C) \  } 0 \qquad \mathrm{(D) \  } 1024 \qquad \mathrm{(E) \  }1024i$


 
 \textbf{Answer: }C



\textbf{Solution}

Notice that $(1+i)^2=2i$ and $(1-i)^2=-2i$. Therefore, 

 
 \[(1+i)^{20}-(1-i)^{20}=(2i)^{10}-(-2i)^{10}=(2i)^{10}-(2i)^{10}=0, \boxed{\text{C}}.\]

 


	\item (\textit{AHSME 1974}) If $\log_8{3}=p$ and $\log_3{5}=q$, then, in terms of $p$ and $q$, $\log_{10}{5}$ equals


 $\mathrm{(A)\ } pq \qquad \mathrm{(B) \ }\frac{3p+q}{5} \qquad \mathrm{(C) \  } \frac{1+3pq}{p+q} \qquad \mathrm{(D) \  } \frac{3pq}{1+3pq} \qquad \mathrm{(E) \  }p^2+q^2$


 
 \textbf{Answer: }D



\textbf{Solution}

Notice that $\frac{1}{p}=\log_{3}{8}$, so it would probably be easier to work in base $3$. From change of base, $\log_{10}{5}=\frac{\log_{3}{5}}{\log_{3}{10}}$. We're given that $\log_{3}{5}=q$, so now we just need to find $\log_{3}{10}$. 

 We have $\frac{1}{p}=\log_{3}{8}=3\log_{3}{2}$, so $\log_{3}{2}=\frac{1}{3p}$. Also, $\log_{3}{10}=\log_{3}{5}+\log{3}{2}=q+\frac{1}{3p}$. Therefore, $\log_{10}{5}=\frac{q}{q+\frac{1}{3p}}=\frac{3pq}{3pq+1}, \boxed{\text{D}}$.

 


	\item (\textit{AHSME 1974}) In the adjoining figure $ABCD$ is a square and $CMN$ is an equilateral triangle. If the area of $ABCD$ is one square inch, then the area of $CMN$ in square inches is 


 
                    \begin{figure}[H]
                        centering
                        \includegraphics[width=0.5\textwidth]{Problem/AHSME_I_1974_19_Problem_Figure_1.png}
                    \end{figure}
                    


 $\mathrm{(A)\ } 2\sqrt{3}-3 \qquad \mathrm{(B) \ }1-\frac{\sqrt{3}}{3} \qquad \mathrm{(C) \  } \frac{\sqrt{3}}{4} \qquad \mathrm{(D) \  } \frac{\sqrt{2}}{3} \qquad \mathrm{(E) \  }4-2\sqrt{3}$


 
 \textbf{Answer: }A



\textbf{Solution}

Let $BN=x$ so that $AN=1-x$. From the Pythagorean Theorem on $\triangle NBC$, we get $CN=\sqrt{x^2+1}$, and from the Pythagorean Theorem on $\triangle AMN$, we get $MN=(1-x)\sqrt{2}$. Since $\triangle CMN$ is equilateral, we must have $\sqrt{x^2+1}=(1-x)\sqrt{2}\implies x^2+1=2x^2-4x+2\implies x^2-4x+1=0$. From the Pythagorean Theorem, we get $x=2-\sqrt{3}$, since we want the root that's less than $1$.

 Therefore, $CN^2=x^2+1=(2-\sqrt{3})^2+1=8-4\sqrt{3}$. The area of an equilateral triangle with side length $x$ is equal to $\frac{x^2\sqrt{3}}{4}$, so the area of $\triangle CMN$ is $\frac{(8-4\sqrt{3})(\sqrt{3})}{4}=2\sqrt{3}-3, \boxed{\text{A}}$.

 



\textbf{Solution 2 (Visualization + Using Ratios)}

We know there is only one way to fit an equilateral triangle into a square: one of its vertices is a corner of the square and the other two vertices fall on opposite sides (try to imagine it in your head). It must be symmetrical along a diagonal of the square.

 Thus $\triangle CNB \sim\triangle CDM$ where both are $15^{\circ}-75^{\circ}-90^{\circ}$ triangles since $\angle MCN=60^{\circ}$

 Using the $15^{\circ}-75^{\circ}-90^{\circ}$ ratios ($2-\sqrt3:1:\sqrt6-\sqrt2$), we have $CN = \sqrt6 - \sqrt2$.

 Thus $[\triangle CMN] = \frac{\sqrt3}{4}(\sqrt6 - \sqrt2)^2 = \frac{\sqrt3}{4}(8-4\sqrt3) = 2\sqrt3-3 \Rightarrow \fbox{A}$.

 $\sim$ Jonysun

 


	\item (\textit{AHSME 1974}) Let


 $T=\frac{1}{3-\sqrt{8}}-\frac{1}{\sqrt{8}-\sqrt{7}}+\frac{1}{\sqrt{7}-\sqrt{6}}-\frac{1}{\sqrt{6}-\sqrt{5}}+\frac{1}{\sqrt{5}-2}.$


 Then


 $\mathrm{(A)\ } T<1 \qquad \mathrm{(B) \ }T=1 \qquad \mathrm{(C) \  } 1<T<2 \qquad \mathrm{(D) \  } T>2 \qquad$


 $\mathrm{(E) \  }T=\frac{1}{(3-\sqrt{8})(\sqrt{8}-\sqrt{7})(\sqrt{7}-\sqrt{6})(\sqrt{6}-\sqrt{5})(\sqrt{5}-2)}$


 
 \textbf{Answer: }D



\textbf{Solution}

Let's try to rationalize $\frac{1}{\sqrt{n+1}-\sqrt{n}}$. Multiplying the numerator and denominator by $\sqrt{n+1}+\sqrt{n}$ gives us $\frac{1}{\sqrt{n+1}-\sqrt{n}}=\sqrt{n+1}+\sqrt{n}$.

 Therefore,

 
 \[T=(3+\sqrt{8})-(\sqrt{8}+\sqrt{7})+(\sqrt{7}+\sqrt{6})-(\sqrt{6}+\sqrt{5})+(\sqrt{5}+2)=5.\]

 Hence the answer is $\boxed{\text{D}}$.

 


	\item (\textit{AHSME 1974}) In a geometric series of positive terms the difference between the fifth and fourth terms is $576$, and the difference between the second and first terms is $9$. What is the sum of the first five terms of this series?


 $\mathrm{(A)\ } 1061 \qquad \mathrm{(B) \ }1023 \qquad \mathrm{(C) \  } 1024 \qquad \mathrm{(D) \  } 768 \qquad \mathrm{(E) \  }\text{none of these}$


 
 \textbf{Answer: }B



\textbf{Solution}

Let the first term be $a$ and the common ratio be $r$. Therefore, the second term is $ar$, the fourth term is $ar^3$, and the fifth term is $ar^4$. We're given that $ar^4-ar^3=576\implies ar^3(r-1)=576$ and $ar-a=9\implies a(r-1)=9$. Dividing this first equation by this second one, we get $r^3=\frac{576}{9}=64\implies r=4$. Therefore, $a(4-1)=9$, so $a=3$. 

 Therefore, the first five terms of this series are $3, 12, 48, 192, 768$, and their sum is $1023, \boxed{\text{B}}$.

 


	\item (\textit{AHSME 1974}) The minimum of $\sin\frac{A}{2}-\sqrt{3}\cos\frac{A}{2}$ is attained when $A$ is 


 $\mathrm{(A)\ } -180^\circ \qquad \mathrm{(B) \ }60^\circ \qquad \mathrm{(C) \  } 120^\circ \qquad \mathrm{(D) \  } 0^\circ \qquad \mathrm{(E) \  }\text{none of these}$


 
 \textbf{Answer: }E



\textbf{Solution}

Define a new number $B$ such that $\cos(B)=\frac{1}{2}$ and $\sin(B)=-\frac{\sqrt{3}}{2}$. Notice that $B=\frac{5\pi}{3}$. Therefore, we have $\sin\frac{A}{2}-\sqrt{3}\cos\frac{A}{2}= 2\cos(B)\sin\left(\frac{A}{2}\right)+2\sin(B)\cos\left(\frac{A}{2}\right)$. 

 Recognizing the angle sum formula, we can rewrite this as $\sin\frac{A}{2}-\sqrt{3}\cos\frac{A}{2}=2\sin\left(\frac{A}{2}+B\right)$. We can now clearly see that the minimum is $-2$, and this is achieved when $\frac{A}{2}+B=\frac{3\pi}{2}$. Plugging in $B=\frac{5\pi}{3}$, we get $\frac{A}{2}=-\frac{\pi}{6}\implies A=-\frac{\pi}{3}=-60^\circ$. $\boxed{\text{E}}$

 


	\item (\textit{AHSME 1974}) In the adjoining figure $TP$ and $T'Q$ are parallel tangents to a circle of radius $r$, with $T$ and $T'$ the points of tangency. $PT''Q$ is a third tangent with $T'''$ as a point of tangency. If $TP=4$ and $T'Q=9$ then $r$ is


 
                    \begin{figure}[H]
                        centering
                        \includegraphics[width=0.5\textwidth]{Problem/AHSME_I_1974_23_Problem_Figure_1.png}
                    \end{figure}
                    


 $\mathrm{(A)\ } 25/6 \qquad \mathrm{(B) \ } 6 \qquad \mathrm{(C) \  } 25/4 \qquad$


 $\mathrm{(D) \  } \text{a number other than }25/6, 6, 25/4 \qquad$


 $\mathrm{(E) \  }\text{not determinable from the given information}$


 
 \textbf{Answer: }B



\textbf{Solution}


                    \begin{figure}[H]
                        centering
                        \includegraphics[width=0.5\textwidth]{Solution/AHSME_I_1974_23_Solution_Figure_1.png}
                    \end{figure}
                    

 Drop the perpendicular from $P$ to $T'Q$ and let the foot be $R$. Note that $PTT'R$ is a rectangle. Also, from Two Tangents, $PT''=4$ and $QT''=9$, so $PQ=13$. Therefore, from the Pythagorean Theorem on $\triangle PRQ$, $PR=\sqrt{13^2-5^2}=12$. We now see that $r=\frac{PR}{2}=6, \boxed{\text{B}}$.

 


	\item (\textit{AHSME 1974}) A fair die is rolled six times. The probability of rolling at least a five at least five times is 


 $\mathrm{(A)\ } \frac{13}{729} \qquad \mathrm{(B) \ }\frac{12}{729} \qquad \mathrm{(C) \  } \frac{2}{729} \qquad \mathrm{(D) \  } \frac{3}{729} \qquad \mathrm{(E) \  }\text{none of these}$


 
 \textbf{Answer: }A



\textbf{Solution}

The probability of rolling at least a five on any one roll is $\frac{2}{6}=\frac{1}{3}$. If there are exactly $5$ fives or sixes rolled, there are $\binom{6}{5}=6$ ways to pick which of the rolls are the fives and sixes, and so the probability in this case is $6\left(\frac{1}{3}\right)^5\left(\frac{2}{3}\right)=\frac{12}{729}$. If there are exactly $6$ fives or sixes rolled, then there is only one way to pick which of the rolls are fives and sixes, so the probability in this case is $\left(\frac{1}{3}\right)^6=\frac{1}{729}$.

 Therefore, the total probability is $\frac{12}{729}+\frac{1}{729}=\frac{13}{729}, \boxed{\text{A}}$.

 


	\item (\textit{AHSME 1974}) In parallelogram $ABCD$ of the accompanying diagram, line $DP$ is drawn bisecting $BC$ at $N$ and meeting $AB$ (extended) at $P$. From vertex $C$, line $CQ$ is drawn bisecting side $AD$ at $M$ and meeting $AB$ (extended) at $Q$. Lines $DP$ and $CQ$ meet at $O$. If the area of parallelogram $ABCD$ is $k$, then the area of the triangle $QPO$ is equal to 


 
                    \begin{figure}[H]
                        centering
                        \includegraphics[width=0.5\textwidth]{Problem/AHSME_I_1974_25_Problem_Figure_1.png}
                    \end{figure}
                    


 $\mathrm{(A)\ } k \qquad \mathrm{(B) \ }\frac{6k}{5} \qquad \mathrm{(C) \  } \frac{9k}{8} \qquad \mathrm{(D) \  } \frac{5k}{4} \qquad \mathrm{(E) \  }2k$


 
 \textbf{Answer: }C



\textbf{Solution}

Note that 
 \[[QPO]=[QAM]+[PBN]+[AMONB]=[AMONB]+[MDC]+[NCD]\] 

 
 \[=[AMONB]+[MDC]+[NOC]+[DOC]=[ABCD]+[DOC]=k+[DOC].\]

 Also, note that $DCNM$ is a parallelogram, and so $[DOC]=\frac{1}{4}[DCNM]=\frac{1}{8}[ABCD]=\frac{k}{8}$.

 Therefore, $[QPO]=k+\frac{k}{8}=\frac{9k}{8}, \boxed{\text{C}}$.

 


	\item (\textit{AHSME 1974}) The number of distinct positive integral divisors of $(30)^4$ excluding $1$ and $(30)^4$ is 


 $\mathrm{(A)\ } 100 \qquad \mathrm{(B) \ }125 \qquad \mathrm{(C) \  } 123 \qquad \mathrm{(D) \  } 30 \qquad \mathrm{(E) \  }\text{none of these}$


 
 \textbf{Answer: }C



\textbf{Solution}

The prime factorization of $30$ is $2\cdot3\cdot5$, so the prime factorization of $30^4$ is $2^4\cdot3^4\cdot5^4$. Therefore, the number of positive divisors of $30^4$ is $(4+1)(4+1)(4+1)=125$. However, we have to subtract $2$ to account for $1$ and $30^4$, so our final answer is $ 125-2=123, \boxed{\text{C}} 

 
soln by RNVAA

 


	\item (\textit{AHSME 1974}) If $f(x)=3x+2$ for all real $x$, then the statement:
"$|f(x)+4|<a$ whenever $|x+2|<b$ and $a>0$ and $b>0$"
is true when


 $\mathrm{(A)}\ b\le a/3\qquad\mathrm{(B)}\ b > a/3\qquad\mathrm{(C)}\ a\le b/3\qquad\mathrm{(D)}\ a > b/3\\ \qquad\mathrm{(E)}\ \text{The statement is never true.}$


 
 \textbf{Answer: }A



\textbf{Solution}

Plugging in $f(x)=3x+2$ into $|f(x)+4|<a$, we get $|3x+6|=|3(x+2)|=3|x+2|<a$. Therefore, $|x+2|<\frac{a}{3}$. Therefore, we want $|x+2|<b\implies |x+2|<\frac{a}{3}$. This is clearly only true when $b\le\frac{a}{3}, \boxed{\text{A}}$.

 


	\item (\textit{AHSME 1974}) Which of the following is satisfied by all numbers $x$ of the form


 
 \[x=\frac{a_1}{3}+\frac{a_2}{3^2}+\cdots+\frac{a_{25}}{3^{25}}\]


 where $a_1$ is $0$ or $2$, $a_2$ is $0$ or $2$,...,$a_{25}$ is $0$ or $2$?


 $\mathrm{(A)\ } 0\le x<1/3 \qquad \mathrm{(B) \ } 1/3\le x<2/3 \qquad \mathrm{(C) \  } 2/3\le x<1 \qquad$


 $\mathrm{(D) \  } 0\le x<1/3 \text{ or }2/3\le x<1 \qquad \mathrm{(E) \  }1/2\le x\le 3/4$


 
 \textbf{Answer: }D



\textbf{Solution}

We can separate this into two cases.

 \textbf{Case 1: $a_1=0$} Here, clearly the minimum sum is when all of the $a_i$ are $0$, so $0\le x$. The maximum sum is when the rest of the $a_i$ are all $2$. In this case, 
 \[\sum_{n=2}^{25}\frac{2}{3^n}<2\sum_{n=2}^{\infty}\frac{1}{3^n}=2\left(\frac{\frac{1}{9}}{1-\frac{1}{3}}\right)=\frac{1}{3}.\] Therefore, in this case, $0\le x<1/3$.

 
\textbf{Case 2: $a_1=2$} Here, again, the minimum sum is when the rest of the $a_i$ are all $0$, in which case the minimum sum is $\frac{2}{3}$. The maximum sum is attained when all of the $a_i$ are $2$. In that case, 
 \[\sum_{n=1}^{25}\frac{2}{3^n}<2\sum_{n=1}^{\infty}\frac{1}{3^n}=2\left(\frac{\frac{1}{3}}{1-\frac{1}{3}}\right)=1.\] Therefore, in this case, $2/3\le x<1$.

 Thus, we must have $0\le x<1/3$ or $2/3\le x<1$, $\boxed{\text{D}}$.

 


	\item (\textit{AHSME 1974}) For $p=1, 2, \cdots, 10$ let $S_p$ be the sum of the first $40$ terms of the arithmetic progression whose first term is $p$ and whose common difference is $2p-1$; then $S_1+S_2+\cdots+S_{10}$ is


 $\mathrm{(A)\ } 80000 \qquad \mathrm{(B) \ }80200 \qquad \mathrm{(C) \  } 80400 \qquad \mathrm{(D) \  } 80600 \qquad \mathrm{(E) \  }80800$


 
 \textbf{Answer: }B



\textbf{Solution}

The $40\text{th}$ term of an arithmetic progression with a first term $p$ and a common difference $2p-1$ is $p+39(2p-1)=79p-39$. Therefore, the sum of the first $40$ terms of such a progression is $\frac{40}{2}(79p-39+p)=1600p-780$.

 We now want to evaluate $\sum_{p=1}^{10}(1600p-780)$. 
 \[\sum_{p=1}^{10}(1600p-780)=1600\sum_{p=1}^{10}(p)-\sum_{p=1}^{10}(780)\] 
 \[=(1600)\left(\frac{10\cdot11}{2}\right)-(780)(10)=88000-7800=80200, \boxed{\text{B}}.\]

 


	\item (\textit{AHSME 1974}) A line segment is divided so that the lesser part is to the greater part as the greater part is to the whole. If $R$ is the ratio of the lesser part to the greater part, then the value of 


 
 \[R^{[R^{(R^2+R^{-1})}+R^{-1}]}+R^{-1}\]


 is


 $\mathrm{(A)\ } 2 \qquad \mathrm{(B) \ }2R \qquad \mathrm{(C) \  } R^{-1} \qquad \mathrm{(D) \  } 2+R^{-1} \qquad \mathrm{(E) \  }2+R$


 
 \textbf{Answer: }A



\textbf{Solution}

Let $w$ be the length of the shorter segment and $l$ be the length of the longer segment. We're given that $\frac{w}{l}=\frac{l}{w+l}$. Cross-multiplying, we find that $w^2+wl=l^2\implies w^2+wl-l^2=0$. Now we divide both sides by $l^2$ to get $\left(\frac{w}{l}\right)^2+\left(\frac{w}{l}\right)-1=0$. Therefore, $R^2+R-1=0$.

 From this, we have $R^2=-R+1$. Dividing both sides by $R$, we get $R=-1+\frac{1}{R}\implies R^{-1}=R+1$. Therefore, $R^2+R^{-1}=-R+1+R+1=2$. Finally, we have 
 \[R^{[R^{(R^2+R^{-1})}+R^{-1}]}+R^{-1}=R^{[R^2+R^{-1}]}+R^{-1}=R^2+R^{-1}= \boxed{\textbf{(A)}2}.\]

 


\end{enumerate}